\documentclass[aspectratio=169, 10pt]{beamer}

\usepackage[utf8]{inputenc}
\usepackage[T1]{fontenc}
\usepackage{xcolor}
\usepackage{graphicx}

\IfFileExists{opensans.sty}{\usepackage[default,scale=0.95]{opensans}}{}

\definecolor{amritaMaroon}{HTML}{A4123F}

\usetheme{Madrid}
\usecolortheme[named=amritaMaroon]{structure}
\setbeamertemplate{navigation symbols}{}
\setbeamertemplate{footline}[frame number]

\begin{document}

% --- SLIDE 1: Technical Change Overview ---
\begin{frame}{Technical Change: Fixed Strategies to Evolved Strategies}

\textbf{\textcolor{amritaMaroon}{Before (Review 1)}}
\begin{itemize}
    \item The algorithm chose from \textbf{14 predefined strategies} like LightRush, HeavyDefense, WorkerRush.
    \item Each strategy was a \textbf{hardcoded script} written by a human programmer.
    \item The AI could only \textbf{select} a strategy --- it could not create one.
    \item This is a \textbf{discrete choice problem} with only 14 possible solutions.
\end{itemize}

\vspace{0.4cm}
\textbf{\textcolor{amritaMaroon}{After (Review 2)}}
\begin{itemize}
    \item The AI now has \textbf{20 continuous weight parameters} (each 0.0 to 1.0).
    \item These weights control every aspect of game behavior --- the AI \textbf{creates its own strategy}.
    \item The search space went from \textbf{14 options to infinite combinations}.
\end{itemize}

\end{frame}


% --- SLIDE 2: The 20 Weights (compact) ---
\begin{frame}{The 20 Evolved Parameters}

{\small
\begin{columns}[T]

\begin{column}{0.32\textwidth}
\textbf{\textcolor{amritaMaroon}{Economy}}\\[0.1cm]
\texttt{w0} worker\_ratio\\
\texttt{w1} max\_workers\\
\texttt{w6} barracks\_timing\\
\texttt{w2} expand\_threshold\\
\texttt{w17} second\_barracks
\end{column}

\begin{column}{0.32\textwidth}
\textbf{\textcolor{amritaMaroon}{Unit Production}}\\[0.1cm]
\texttt{w3} light\_preference\\
\texttt{w4} heavy\_preference\\
\texttt{w5} ranged\_preference\\
\texttt{w13} tech\_switch
\end{column}

\begin{column}{0.32\textwidth}
\textbf{\textcolor{amritaMaroon}{Combat \& Defense}}\\[0.1cm]
\texttt{w7} aggression\\
\texttt{w8} attack\_threshold\\
\texttt{w9} defend\_radius\\
\texttt{w10} target\_priority\\
\texttt{w11} defense\_trigger\\
\texttt{w12} counter\_rush\\
\texttt{w14} economy\_pressure\\
\texttt{w15} scout\_priority\\
\texttt{w16} group\_attack\\
\texttt{w18} worker\_fight\_ratio\\
\texttt{w19} focus\_fire
\end{column}

\end{columns}
}

\vspace{0.4cm}
Each weight is a continuous value between 0.0 and 1.0. The combination of all 20 weights defines a \textbf{unique strategy} that no built-in bot has.

\end{frame}


% --- SLIDE 3: New Approaches ---
\begin{frame}{New Approaches for Improving TD-CCEA}

\textbf{\textcolor{amritaMaroon}{Approach 1: Sensitivity Analysis (SA-TDCCEA)}}
\begin{itemize}
    \item Not all 20 weights are equally important. Some heavily affect the outcome, others barely matter.
    \item We measure each weight's \textbf{influence} on the game result independently.
    \item Important weights get \textbf{careful, small mutations}. Unimportant weights get \textbf{large, free mutations}.
    \item Focuses the evolution effort where it matters most.
\end{itemize}

\vspace{0.4cm}
\textbf{\textcolor{amritaMaroon}{Approach 2: Dynamic Programming (DP-TDCCEA)}}
\begin{itemize}
    \item The same opponents repeat every 7 generations --- this creates \textbf{overlapping subproblems}.
    \item Instead of solving each generation from scratch, we \textbf{store and reuse} past results.
    \item The trajectory memory $H_i$ acts as a \textbf{DP table}, storing (fitness, environment) pairs.
    \item Agents are scored based on \textbf{consistent performance} across similar past situations.
\end{itemize}

\end{frame}


% --- SLIDE 4: Sensitivity Analysis ---
\begin{frame}{Sensitivity Analysis: How It Works}

\textbf{\textcolor{amritaMaroon}{Step 1: Measure Influence}}
\begin{itemize}
    \item Set all 20 weights to 0.5 (neutral baseline).
    \item For each weight individually:
    \begin{itemize}
        \item Set it to 0.0 (all others stay 0.5) $\rightarrow$ play games $\rightarrow$ measure win rate.
        \item Set it to 1.0 (all others stay 0.5) $\rightarrow$ play games $\rightarrow$ measure win rate.
    \end{itemize}
    \item \textbf{Influence} = $|$win rate at 1.0 $-$ win rate at 0.0$|$
\end{itemize}

\vspace{0.3cm}
\textbf{\textcolor{amritaMaroon}{Step 2: Assign Per-Weight Mutation Rates}}
\begin{itemize}
    \item Standard TD-CCEA: $\mu = 0.15$ for all 20 weights (same everywhere).
    \item SA-TDCCEA: each weight gets its \textbf{own mutation rate} based on influence.
    \begin{itemize}
        \item High influence (e.g., \texttt{aggression}) $\rightarrow$ $\mu_j = 0.08$ (fine-tune carefully).
        \item Low influence (e.g., \texttt{second\_barracks}) $\rightarrow$ $\mu_j = 0.20$ (explore freely).
    \end{itemize}
    \item Like studying harder for subjects that carry more marks.
\end{itemize}

\end{frame}


% --- SLIDE 5: Why DP Is Good ---
\begin{frame}{Dynamic Programming: Why It Works for TD-CCEA}

\textbf{\textcolor{amritaMaroon}{The Problem It Solves}}
\begin{itemize}
    \item Opponents rotate every 7 generations: the \textbf{same opponent reappears} multiple times.
    \item Without DP, the algorithm evaluates agents from scratch each time --- even if it faced the same situation before.
    \item This wastes computation and \textbf{forgets} what strategies worked in the past.
\end{itemize}

\vspace{0.3cm}
\textbf{\textcolor{amritaMaroon}{How DP Helps}}
\begin{itemize}
    \item Each agent stores its last 7 results as (fitness, environment) pairs in a \textbf{trajectory memory}.
    \item When the same or similar environment appears again, the agent \textbf{already knows} how well it performed.
    \item The agent's score is a \textbf{weighted average} of all past results --- not just the latest game.
    \item Past results from similar situations are \textbf{trusted more}, results from different situations are \textbf{trusted less}.
\end{itemize}

\vspace{0.3cm}
\textbf{\textcolor{amritaMaroon}{The Benefit}}
\begin{itemize}
    \item An agent that won once by luck gets \textbf{exposed} --- its average across all past situations is low.
    \item An agent that wins consistently across similar opponents \textbf{rises to the top}.
    \item The algorithm avoids re-solving problems it has already solved --- \textbf{classic DP principle}.
\end{itemize}

\end{frame}


% --- SLIDE 6: Why DP Is Better ---
\begin{frame}{Why DP Proves to Be Better}

\textbf{\textcolor{amritaMaroon}{Without DP (Greedy)}}
\begin{itemize}
    \item An agent wins one game by luck $\rightarrow$ gets high fitness $\rightarrow$ survives.
    \item Next generation against a different opponent $\rightarrow$ loses badly.
    \item No way to tell a \textbf{lucky agent} from a \textbf{genuinely good agent}.
    \item The algorithm is easily fooled by single-generation noise.
\end{itemize}

\vspace{0.4cm}
\textbf{\textcolor{amritaMaroon}{With DP (Trajectory Memory)}}
\begin{itemize}
    \item The trajectory score $e_i$ averages performance across \textbf{all similar past situations}.
    \item Lucky one-time winners get \textbf{diluted} --- their single high score is averaged with past losses.
    \item Consistently good agents \textbf{rise to the top} because they score well across multiple environments.
    \item Adaptive mutation is also based on $e_i$: good agents are \textbf{fine-tuned}, bad agents \textbf{explore more}.
    \item The algorithm makes \textbf{smarter survival decisions} based on accumulated evidence, not just the latest result.
\end{itemize}

\end{frame}


\end{document}
